
\documentclass[11pt,a4paper]{article}

\usepackage[utf8]{inputenc}
\usepackage[T1]{fontenc}
\usepackage[spanish,es-noquoting]{babel}
\usepackage{lmodern}
\usepackage{geometry}
\usepackage{microtype}
\usepackage{graphicx}
\usepackage{array}
\usepackage{booktabs}
\usepackage{longtable}
\usepackage{tabularx}
\usepackage{enumitem}
\usepackage{tikz}
\usepackage{float}
\usepackage{fancyhdr}
\usepackage{titlesec}
\usepackage[hidelinks]{hyperref}
\usepackage{xurl}
\usepackage{url}
\usetikzlibrary{arrows.meta,positioning,fit,calc}

\geometry{left=2.6cm,right=2.6cm,top=2.6cm,bottom=2.5cm}
\setlength{\parindent}{0pt}
\setlength{\parskip}{6pt}
\emergencystretch=3em
\setcounter{tocdepth}{2}
\setlist[itemize]{leftmargin=*,topsep=2pt,itemsep=2pt}
\setlist[enumerate]{leftmargin=*,topsep=2pt,itemsep=2pt}

\pagestyle{fancy}
\fancyhf{}
\lhead{\small Arquitectura de Software}
\rhead{\small Laboratorio 2 -- PersonAPP}
\cfoot{\thepage}
\renewcommand{\headrulewidth}{0.3pt}
\setlength{\headheight}{15pt}
\titleformat{\section}{\large\bfseries}{\thesection.}{0.6em}{}
\titleformat{\subsection}{\normalsize\bfseries}{\thesubsection.}{0.6em}{}

\newcommand{\pendiente}[1]{\textbf{Pendiente de completar:} #1}
\newcommand{\base}{\textit{plantilla del profesor}}

\begin{document}
\begin{titlepage}
\thispagestyle{empty}
\begin{center}
{\large Pontificia Universidad Javeriana\par}
\vspace{0.25cm}
{Facultad de Ingeniería\par}
{Departamento de Ingeniería de Sistemas\par}
\vspace{3.0cm}
{\LARGE\bfseries Laboratorio 2\par}
\vspace{0.45cm}
{\Large Implementación de un servicio web\par}
{\Large con arquitectura hexagonal, Repository y Service\par}
\vspace{0.65cm}
{\large PersonAPP -- Spring Boot\par}
\vspace{0.8cm}
{JDK 11 \(\cdot\) Spring Boot \(\cdot\) Maven\par}
{MariaDB y MongoDB \(\cdot\) REST y CLI \(\cdot\) Swagger 3\par}
\vfill
\textbf{Integrantes}\par
Sara Albarracín\par
Juan Diego Rojas\par
Miguel Sarmiento\par
\vspace{1cm}
\textbf{Profesor}\par
Andrés Armando Sánchez Martín\par
\vspace{1cm}
Arquitectura de Software\par
Bogotá D. C., noviembre de 2026\par
\end{center}
\end{titlepage}
\tableofcontents
\clearpage

\section{Marco conceptual}
El Laboratorio 2 consiste en desarrollar un servicio web con operaciones CRUD
(\emph{Create, Read, Update, Delete}) para un modelo de personas, profesiones,
estudios y teléfonos. A diferencia del Laboratorio 1, cuyo diseño se centraba
en un monolito .NET con MVC y DAO, en este ejercicio se emplean Java, Spring Boot,
Maven y la organización por puertos y adaptadores propia de la arquitectura hexagonal.
La solución debe contemplar dos clientes (REST y CLI), dos motores de persistencia
(MariaDB y MongoDB) y documentación de los servicios mediante Swagger 3 (OpenAPI 3).
El punto de partida es el repositorio de referencia proporcionado por el profesor
\cite{plantilla}.

El stack del laboratorio es Java, Spring Boot, MongoDB y MariaDB, REST y CLI,
y Swagger 3. La plantilla declara JDK 11 y Spring Boot 2.7.11, que son la línea
base. Según la explicación del profesor, después de hacerla funcionar se
actualizan sus tecnologías a versiones recientes y compatibles, se contenedoriza
la aplicación y, al final, se completan las funcionalidades.

El equipo desglosó el CRUD de Persona, Profesión, Teléfono y Estudios en seis
operaciones: tres lecturas (\textbf{consultar por identificador}, \textbf{listar
todos} y \textbf{contar} los registros del motor seleccionado) y tres escrituras
(\textbf{crear}, \textbf{editar} y \textbf{eliminar}). El listado parcial con
paginación y orden estable se conserva como operación adicional. También es decisión del equipo permitir el
cambio de motor durante la ejecución, sin reiniciar la aplicación. Ambos son
criterios de integración del equipo, no requisitos textuales del enunciado.

\subsection{Arquitectura de software}
La arquitectura de software describe las estructuras de un sistema, las
responsabilidades de sus elementos y sus relaciones. Una decisión relevante en
este laboratorio es aislar la lógica del negocio de las interfaces de usuario
y de los mecanismos de persistencia; esto facilita la sustitución de adaptadores,
las pruebas y el mantenimiento.

\subsection{Arquitectura limpia y arquitectura hexagonal}
La arquitectura hexagonal, también denominada \emph{Ports and Adapters},
organiza el sistema alrededor del núcleo de negocio, que se comunica mediante
contratos explícitos \cite{cockburn}. Los \textbf{puertos de entrada} describen
las operaciones que pueden solicitar los clientes y los \textbf{puertos de salida}
definen los servicios externos que requieren los casos de uso, especialmente
la persistencia. Los \textbf{adaptadores} conectan estos contratos con
tecnologías concretas.

La arquitectura limpia comparte la idea de orientar las dependencias hacia
la lógica central y evitar que reglas de negocio dependan de frameworks o bases
de datos concretas \cite{martin}. En este laboratorio se expresa mediante
módulos Maven: \texttt{domain} contiene los modelos de negocio;
\texttt{application} incorpora los puertos y casos de uso;
\texttt{rest-input-adapter} y \texttt{cli-input-adapter} reciben solicitudes;
\texttt{maria-output-adapter} y \texttt{mongo-output-adapter} realizan la
persistencia. El módulo \texttt{common} reúne elementos compartidos. La existencia
de estos módulos expresa el diseño previsto, aunque la dirección real de las
dependencias debe revisarse en el código.

\subsection{Modularidad y Maven}
La modularidad divide la aplicación en unidades con responsabilidades específicas.
Maven define su compilación, dependencias y configuración mediante archivos
\texttt{pom.xml}. La plantilla es un proyecto multimódulo con un POM padre y
siete módulos, cada uno con su propio POM \cite{maven,plantilla}. Esta división
facilita localizar los puntos de cambio, pero no implica que cada módulo sea
un microservicio independiente.

\subsection{Puertos, adaptadores, casos de uso, Service y Repository}
El patrón \textbf{Repository} encapsula la persistencia de las entidades y
proporciona operaciones de consulta y modificación. En la plantilla,
\texttt{PersonOutputPort} define un contrato con \texttt{save},
\texttt{delete}, \texttt{find} y \texttt{findById}; los adaptadores de salida
implementan dicho contrato utilizando los repositorios de cada motor.

Un \textbf{Service} o servicio de aplicación coordina operaciones del sistema.
El caso de uso \texttt{PersonUseCase} implementa \texttt{PersonInputPort} y
contiene operaciones como \texttt{create}, \texttt{edit}, \texttt{drop},
\texttt{findAll}, \texttt{findOne}, \texttt{count}, \texttt{getPhones} y
\texttt{getStudies}. Este caso de uso no debe confundirse con el controlador REST:
un método definido en la aplicación no necesariamente está expuesto todavía
por HTTP o por el menú de consola.

\subsection{Inyección de dependencias y selección de persistencia}
Spring dispone de un contenedor para crear e inyectar objetos. Cuando dos
implementaciones cumplen la misma interfaz, una selección explícita, por
ejemplo mediante \texttt{@Qualifier}, permite identificar el componente
requerido \cite{spring}. La plantilla define adaptadores de salida de MariaDB
y MongoDB que implementan \texttt{PersonOutputPort}; los adaptadores de entrada
seleccionan la implementación según una opción como \texttt{MARIA} o
\texttt{MONGO}. El equipo se propuso conservar esa capacidad de cambio durante
la ejecución y probarla en los dos clientes.

La explicación del profesor asocia esta posibilidad con el concepto de CQRS,
pero \textbf{elegir una base u otra no implica automáticamente implementar CQRS}:
ese patrón separa las responsabilidades de lectura y escritura, lo que requiere
un diseño más específico \cite{fowler}.

\subsection{DTO y Mapper}
Un DTO (\emph{Data Transfer Object}) transporta la información entre una
interfaz externa y la aplicación. Un Mapper convierte representaciones de
datos entre capas, por ejemplo, una entidad de dominio \texttt{Person} y una
entidad de MariaDB \texttt{PersonaEntity}, o un documento de MongoDB
\texttt{PersonaDocument}. En la plantilla aparecen los métodos de conversión
\texttt{fromAdapterToDomain} y \texttt{fromDomainToAdapter}, y sus equivalentes
específicos de cada adaptador. El profesor advierte que existe un error en un
Mapper: debe localizarse mediante ejecución y diagnóstico, sin asumir cuál
es el fallo antes de reproducirlo.

\subsection{Spring Boot, Java y las versiones del proyecto}
Java es el lenguaje de implementación; el JDK proporciona las herramientas
necesarias para compilar y ejecutar el programa. Spring Boot permite configurar
y ejecutar aplicaciones Java y servicios HTTP \cite{spring}. El repositorio
base declara \textbf{Java 11} y \textbf{Spring Boot 2.7.11} en el POM principal.
El adaptador REST incorpora \texttt{springdoc-openapi-ui} 1.7.0, que genera la
documentación OpenAPI 3 y la interfaz Swagger UI. Estas versiones son la línea
base con la que se repara la plantilla. Después se actualizan a la versión LTS
más alta compatible donde exista esa clasificación y, para el resto, a la versión
estable más reciente con soporte vigente (HU-007). Actualizar no se reduce a
cambiar números en el POM: obliga a revisar las dependencias directas y
transitivas, los plugins, las anotaciones y las API que cambian entre versiones.
Las versiones finales se registrarán cuando la migración esté verificada.

\subsection{MariaDB y MongoDB}
MariaDB es un sistema de bases de datos relacional que representa los datos
en tablas, índices, claves primarias y relaciones mediante claves foráneas.
MongoDB es un sistema orientado a documentos, por lo que las entidades pueden
guardarse como documentos con estructuras propias \cite{mariadb,mongodb}.
En esta práctica se usan los adaptadores de salida correspondientes, sin
acoplar el caso de uso a un único sistema gestor.

\subsection{REST, CLI, Swagger y operaciones CRUD}
REST es un estilo de interacción basado en recursos y en la semántica de HTTP.
Para un CRUD se emplean habitualmente \texttt{GET} para consultar,
\texttt{POST} para crear, \texttt{PUT}/\texttt{PATCH} para actualizar y
\texttt{DELETE} para eliminar. CLI (\emph{Command Line Interface}) permite
interactuar desde una terminal. Swagger UI proporciona una interfaz gráfica
para consultar y probar una API descrita mediante OpenAPI \cite{openapi}.

\subsection{Scripts DDL y DML}
DDL (\emph{Data Definition Language}) agrupa instrucciones para definir
estructuras, mientras que DML (\emph{Data Manipulation Language}) comprende
operaciones para modificar los datos. El repositorio incluye los archivos
\path{persona_ddl_maria.sql}, \path{persona_dml_maria.sql},
\path{persona_ddl_mongo.js} y \path{persona_dml_mongo.js}.
La estructura física de una base relacional no debe confundirse con la
representación documental de MongoDB.

\section{Diseño}

\subsection{Objetivo y alcance}
Diseñar, completar y documentar una aplicación Java organizada por módulos
hexagonales para gestionar personas, profesiones, estudios y teléfonos,
con operaciones CRUD a través de REST y CLI, persistencia en MariaDB y MongoDB
y documentación de los endpoints. El trabajo parte de un fork de la plantilla
oficial y debe publicarse con un tag Git y los scripts solicitados.

\subsection{Comparación con el Laboratorio 1}
\begin{table}[H]
\centering\small
\begin{tabularx}{\textwidth}{@{}p{3.1cm}XX@{}}
\toprule
\textbf{Aspecto}&\textbf{Laboratorio 1}&\textbf{Laboratorio 2}\\\midrule
Lenguaje / plataforma& C\# y .NET& Java y Spring Boot\\
Diseño principal& Monolito con MVC y DAO& Arquitectura hexagonal, limpia y modular\\
Acceso a datos& Entity Framework Core y SQL Server& Adaptadores Repository, MariaDB y MongoDB\\
Clientes& Vistas web MVC y API REST& API REST y aplicación CLI\\
Persistencia& Una tecnología central& Selección dinámica de motor\\
Entrega técnica& Repositorio y etiqueta Git& Fork actualizado, scripts, README, código y etiqueta Git\\
\bottomrule
\end{tabularx}
\caption{Cambios principales entre los dos laboratorios.}
\end{table}

\subsection{Requerimientos de la actividad}
\begin{longtable}{@{}p{1.8cm}p{11.9cm}@{}}
\toprule
\textbf{Código}&\textbf{Requerimiento}\\\midrule
\endfirsthead
\toprule\textbf{Código}&\textbf{Requerimiento}\\\midrule
\endhead
RF-01 & Implementar el CRUD de Persona, Profesión, Estudios y Teléfono. El equipo definió tres lecturas (consultar por identificador, listar todos y contar) y tres escrituras (crear, editar y eliminar); el listado parcial es una operación adicional.\\
RF-02 & Publicar los endpoints REST correspondientes y documentarlos en Swagger 3 (OpenAPI 3).\\
RF-03 & Proporcionar operaciones equivalentes mediante un cliente CLI según el alcance definido.\\
RF-04 & Conectar los casos de uso con MariaDB y MongoDB mediante adaptadores de salida.\\
RF-05 & Como criterio del equipo, permitir escoger MariaDB o MongoDB en tiempo de ejecución, sin reiniciar la aplicación.\\
RF-09 & Consultar una persona con todos sus teléfonos por REST y CLI en ambos motores.\\
RA-01 & Organizar las dependencias por módulos, siguiendo el diseño hexagonal y limpio, sin añadir, fusionar ni renombrar los siete módulos de la plantilla.\\
RA-02 & Emplear Repository, Service/casos de uso, puertos y adaptadores.\\
RT-01 & Partir de la plantilla con JDK 11, Spring Boot 2.7.11 y Maven multimódulo, y hacerla funcionar.\\
RT-02 & Actualizar Java, Spring Boot, librerías y plugins a versiones recientes y compatibles, resolviendo dependencias y cambios de API.\\
RT-04 & Contenerizar la aplicación REST y CLI junto con MariaDB y MongoDB.\\
RE-01 & Incluir README con configuración del ambiente, compilación y despliegue.\\
RE-02 & Entregar scripts DDL/DML, código y URL del tag Git del repositorio del equipo.\\
RE-03 & Entregar portada, marco conceptual, diseño, procedimiento, conclusiones, lecciones y referencias.\\
\bottomrule
\caption{Requerimientos sistematizados a partir del enunciado y la explicación del profesor.}
\end{longtable}

\subsection{Modelo de datos}
El modelo funcional es el mismo del Laboratorio 1. La plantilla conserva
las entidades de persona, profesion, estudios y telefono. El archivo
\texttt{scripts/persona\_ddl\_maria.sql} define las tablas que se resumen
a continuación.

\begin{table}[H]
\centering\small
\begin{tabularx}{\textwidth}{@{}p{2.2cm}p{5.2cm}X@{}}
\toprule
\textbf{Tabla}&\textbf{Campos}&\textbf{Clave y relaciones}\\\midrule
\texttt{persona}&\texttt{cc}, \texttt{nombre}, \texttt{apellido}, \texttt{genero}, \texttt{edad}&PK: \texttt{cc}; genero restringido a M/F en el DDL base.\\
\texttt{profesion}&\texttt{id}, \texttt{nom}, \texttt{des}&PK: \texttt{id}.\\
\texttt{estudios}&\texttt{id\_prof}, \texttt{cc\_per}, \texttt{fecha}, \texttt{univer}&PK compuesta por \texttt{id\_prof} y \texttt{cc\_per}. FK hacia profesion y persona.\\
\texttt{telefono}&\texttt{num}, \texttt{oper}, \texttt{duenio}&PK: \texttt{num}; FK \texttt{duenio} hacia persona.\\
\bottomrule
\end{tabularx}
\caption{Modelo relacional del repositorio base en MariaDB.}
\end{table}

Una persona puede asociarse con varios números telefónicos y varios registros
de estudios; cada estudio enlaza una persona con una profesión. En MongoDB
se debe describir cómo se representan estos datos como documentos y verificar
que los adaptadores preserven su significado.

\subsection{Arquitectura propuesta y módulos Maven}
El archivo \texttt{pom.xml} principal de la plantilla declara siete módulos.
La Tabla \ref{tab:modulos} relaciona sus responsabilidades.

\begin{table}[H]
\centering\small
\begin{tabularx}{\textwidth}{@{}p{4.0cm}X@{}}
\toprule
\textbf{Módulo}&\textbf{Responsabilidad}\\\midrule
\texttt{common}&Anotaciones personalizadas, excepciones y opciones de configuración compartidas.\\
\texttt{domain}&Objetos de negocio: Person, Phone, Profession, Study y Gender.\\
\texttt{application}&Puertos de entrada/salida y casos de uso (p. ej., \texttt{PersonUseCase}).\\
\texttt{rest-input-adapter}&Inicio de Spring REST, controlador, DTO y Mapper para HTTP.\\
\texttt{cli-input-adapter}&Inicio de Spring CLI, menús, DTO y Mapper de consola.\\
\texttt{maria-output-adapter}&Entidades JPA, repositorios y mapeadores MariaDB.\\
\texttt{mongo-output-adapter}&Documentos, repositorios y mapeadores MongoDB.\\
\bottomrule
\end{tabularx}
\caption{Estructura del proyecto entregado como plantilla.}
\label{tab:modulos}
\end{table}

\begin{figure}[H]
\centering
\begin{tikzpicture}[
 every node/.style={font=\footnotesize,align=center},
 blk/.style={draw,rounded corners=2pt,minimum height=1cm,text width=3.4cm,fill=white},
 core/.style={draw,very thick,rounded corners=2pt,minimum height=1cm,text width=3.8cm,fill=white},
 arr/.style={-{Latex[length=2mm]},line width=.7pt}
]
\node[blk] (rest) at (-4.3,2.6) {REST / HTTP\\\texttt{rest-input-adapter}};
\node[blk] (cli) at (4.3,2.6) {Terminal / CLI\\\texttt{cli-input-adapter}};
\node[core] (app) at (0,0.65) {Puertos y casos de uso\\\texttt{application}};
\node[core] (dom) at (0,-1.1) {Modelos de negocio\\\texttt{domain}};
\node[blk] (ma) at (-4.3,-2.9) {\texttt{maria-output-adapter}\\MariaDB};
\node[blk] (mo) at (4.3,-2.9) {\texttt{mongo-output-adapter}\\MongoDB};
\draw[arr] (rest) -- (app);
\draw[arr] (cli) -- (app);
\draw[arr] (app) -- (dom);
\draw[arr] (app) -- (ma);
\draw[arr] (app) -- (mo);
\end{tikzpicture}
\caption{Interacciones conceptuales de la arquitectura hexagonal. Las flechas representan el recorrido lógico, no dependencias de código compilado.}
\end{figure}

En la implementación, los adaptadores de salida materializan el contrato del
puerto de persistencia y deben depender de abstracciones definidas en
\texttt{application}; las entidades de \texttt{domain} deben mantenerse
independientes de esos detalles externos. Las dos clases de arranque de
Spring se ubican en los adaptadores de entrada.

\subsection{Flujo de una petición y selección de persistencia}
Como ejemplo, en la plantilla la consulta de personas por REST sigue este
recorrido: \texttt{PersonaController\allowbreak V1} recibe la petición,
\texttt{PersonaInputAdapter\allowbreak Rest} interpreta la base solicitada,
\texttt{PersonUseCase} ejecuta el caso de uso mediante
\texttt{PersonOutputPort}, el adaptador de salida consulta el repositorio
elegido y el Mapper construye la respuesta. La CLI realiza un recorrido
análogo desde el menú de consola.

El equipo decidió que la selección dinámica sea un criterio de diseño: el usuario
o la solicitud indica la base que se usará, y la aplicación elige la implementación
de persistencia correspondiente. Durante las pruebas deberá comprobarse que
distintas peticiones puedan usar ambas bases sin detener el servicio.

\subsection{Diseño de las seis operaciones}
El diseño acordado por el equipo aplica las mismas seis operaciones a Persona,
Profesión, Teléfono y Estudios, en REST y en CLI, sobre el motor seleccionado.
Son tres lecturas y tres escrituras:
\begin{enumerate}
\item \textbf{Consultar por identificador}: devolver un único registro a partir de
su clave completa; si no existe en el motor elegido, se informa un error controlado.
\item \textbf{Listar todos} los registros, ordenados por clave.
\item \textbf{Contar} los registros del motor seleccionado; con una colección vacía
el resultado es cero.
\item \textbf{Crear} un registro.
\item \textbf{Editar} un registro existente identificado por su clave.
\item \textbf{Eliminar} un registro existente.
\end{enumerate}
Como operación adicional, el \textbf{listado parcial} devuelve una página de tamaño
acotado con orden estable por clave, de modo que las páginas sucesivas no repitan
ni omitan registros. Además, HU-033 consulta una persona con todos sus teléfonos.
Estudios usa la clave compuesta (\texttt{id\_prof}, \texttt{cc\_per}) del modelo
base. Las rutas REST, las opciones de menú y los códigos de respuesta están en los
contratos de la wiki del equipo y se describirán aquí cuando estén implementados.

\subsection{Implementación observada en el código actual}
Esta subsección describe el código del repositorio del equipo en el commit
\texttt{a83523b}. Hasta ese commit el equipo solo ha modificado documentación, de
modo que el código coincide con la plantilla del profesor. La tabla distingue lo
declarado de lo efectivamente expuesto. No describe resultados de ejecución.

\begin{table}[H]
\centering\small
\begin{tabularx}{\textwidth}{@{}p{3.6cm}X@{}}
\toprule
\textbf{Elemento}&\textbf{Estado observado en el código}\\\midrule
\texttt{PersonInputPort} y \texttt{PersonUseCase}&Declaran e implementan crear, editar, eliminar, listar, consultar por identificador, contar y consultar teléfonos/estudios. \texttt{count} se calcula como el tamaño de \texttt{findAll()} y no se expone en REST ni en CLI. No existe listado paginado.\\
\texttt{PersonOutputPort}&Ofrece \texttt{save}, \texttt{delete}, \texttt{find} y \texttt{findById}. Lo implementan \texttt{PersonOutputAdapterMaria} y \texttt{PersonOutputAdapterMongo}, que cuentan con sus repositorios.\\
\texttt{PersonaControllerV1}&Declara solo \texttt{GET /api/v1/persona/\{database\}} (listado) y \texttt{POST /api/v1/persona} (alta).\\
\texttt{PersonaMapperRest}&\texttt{fromAdapterToDomain} está incompleto: devuelve un \texttt{Person} vacío (\texttt{TODO}) y no copia los datos de \texttt{PersonaRequest}. La respuesta del alta se etiqueta siempre como MariaDB.\\
\texttt{MenuPrincipal}&Muestra los módulos Persona, Profesión, Teléfono y Estudios. Solo Persona abre un submenú; las otras opciones registran ``Implementar Menu''.\\
\texttt{PersonaMenu}&Permite elegir MariaDB (1) o MongoDB (2) y conecta únicamente la opción 1, que lista todas las personas.\\
Profesión, Teléfono y Estudios&Existen modelos de dominio, entidades JPA, documentos MongoDB y mappers en ambos adaptadores de salida. No existen puertos, casos de uso, repositorios, adaptadores de salida ni rutas o menús.\\
Pruebas automatizadas&No hay directorios \texttt{src/test} en los módulos.\\
\bottomrule
\end{tabularx}
\caption{Implementación observada en el código del commit \texttt{a83523b}, sin ejecución.}
\end{table}

En \texttt{PersonaRequest} los campos son \texttt{dni}, \texttt{firstName},
\texttt{lastName}, \texttt{age}, \texttt{sex} y \texttt{database}. Los
adaptadores de entrada guardan el caso de uso elegido en un atributo del bean,
por lo que el aislamiento entre solicitudes concurrentes que eligen motores
distintos debe comprobarse en las pruebas.

\subsection{Pruebas ejecutadas}
Las pruebas de la fase 1 se ejecutaron en Windows 11 con JDK 11, con MariaDB y
MongoDB en contenedores Docker (\texttt{mariadb:11.4} y \texttt{mongo:6.0}). La
única prueba automatizada es la regresión del Mapper REST de HU-006
(\texttt{PersonaMapperRestTest}). Los comandos completos están en la página de
ejecución local de la wiki del repositorio.

\begin{table}[H]
\centering\small
\begin{tabularx}{\textwidth}{@{}>{\raggedright\arraybackslash}p{2.9cm}>{\raggedright\arraybackslash}X>{\raggedright\arraybackslash}p{2.3cm}@{}}
\toprule
\textbf{Prueba}&\textbf{Resultado observado}&\textbf{Fecha y PR}\\\midrule
Ambiente (HU-003)&Temurin 11.0.32.1, Maven 3.9.9, MariaDB 11.4.13 y MongoDB 6.0.28.&2026-11-08, \#33\\
Compilación&\texttt{mvn -B clean install -DskipTests} termina con \texttt{BUILD SUCCESS} en los ocho artefactos con JDK 11.&2026-11-08, \#33; 2026-11-09, \#35\\
Scripts DDL/DML (HU-004)&MariaDB crea las tablas \texttt{persona}, \texttt{profesion}, \texttt{telefono} y \texttt{estudios} y carga cinco personas. MongoDB carga las mismas cinco personas en \texttt{persona\_db.persona}.&2026-11-09, \#34\\
REST (HU-005)&\texttt{GET /api/v1/persona/MARIA} y \texttt{/MONGO} responden 200 con cinco personas cada uno; un motor desconocido responde 200 con una lista vacía. \texttt{POST /api/v1/persona} responde 500 por una \texttt{NullPointerException} en \texttt{PersonaMapperMaria.validateGenero} y no guarda datos.&2026-11-09, \#35\\
Swagger (HU-005)&\path{/swagger-ui.html} responde 200 y redirige a \path{/swagger-ui/index.html}. \path{/api-docs} publica OpenAPI 3.0.1 con las dos rutas.&2026-11-09, \#35\\
CLI (HU-005)&En una sola sesión lista las cinco personas de MariaDB y luego las de MongoDB, sin reiniciar. Profesión registra ``Implementar Menu''.&2026-11-09, \#35\\
Mapper REST (HU-006)&\texttt{mvn -B clean install} termina con \texttt{BUILD SUCCESS} y \texttt{Tests run: 6, Failures: 0}. \texttt{POST /api/v1/persona} responde 200 y guarda en MariaDB y en MongoDB, también sin edad. Un \texttt{sex} no admitido responde 500 y no guarda datos.&2026-11-10/11, \#38\\
\bottomrule
\end{tabularx}
\caption{Pruebas ejecutadas en la fase 1.}
\end{table}

\subsection{Parámetros iniciales de configuración}
\begin{table}[H]
\centering\small
\begin{tabularx}{\textwidth}{@{}p{3.2cm}p{3.3cm}X@{}}
\toprule
\textbf{Componente}&\textbf{Valor en plantilla}&\textbf{Ubicación o nota}\\\midrule
REST&Puerto 3000&\texttt{application.properties} del adaptador REST\\
OpenAPI 3&\texttt{/api-docs}&\texttt{springdoc.api-docs.path}; \texttt{springdoc-openapi-ui} 1.7.0\\
Swagger UI&\url{http://localhost:3000/swagger-ui.html}&Verificada en HU-005; redirige a \path{/swagger-ui/index.html}\\
MariaDB&Puerto 3307&\texttt{jdbc:mariadb://localhost:3307/persona\_db}\\
MongoDB&Puerto 27017&Base \texttt{persona\_db} en propiedades de Spring\\
Java&11&\texttt{global.java.version} en POM padre\\
Spring Boot&2.7.11&\texttt{global.springframework.boot.version} en POM padre\\
\bottomrule
\end{tabularx}
\caption{Valores de referencia del código que se verificarán en la entrega.}
\end{table}

\section{Procedimiento}

Esta sección documenta la \emph{secuencia prevista} con base en la explicación
del profesor. Los subapartados marcados como pendientes se completarán con
comandos y evidencias \textbf{realmente obtenidos}, no con resultados supuestos.
La guía no trae un tutorial detallado: el grupo debe ejecutar, diagnosticar,
actualizar, contenedorizar y completar el proyecto.

El profesor fijó cuatro tareas en este orden: (1) hacer funcionar la plantilla
original, (2) actualizar sus tecnologías, (3) contenedorizarla y (4) completar
sus funcionalidades. El trabajo se organiza en esas cuatro fases y cada una se
valida antes de iniciar la siguiente. En todas se conservan los siete módulos
Maven, sus puertos y sus adaptadores: el profesor advirtió que pedir una
corrección ambigua puede terminar cambiando la arquitectura, y el objetivo es
aprender a trabajar con módulos y adaptadores. Al final se verifica y se publica
la entrega.

\subsection{Fase 1. Hacer funcionar la plantilla original}
Criterio de salida: el ambiente se prepara con JDK 11, ambos motores son
accesibles, REST y CLI arrancan y el defecto de mapeo tiene una prueba de
regresión. Se corrigen solo los defectos de mapeo, anotaciones, librerías y
dependencias que impidan ejecutar la plantilla.
\subsubsection*{Actividad 1. Preparar el repositorio del equipo}
El profesor solicita crear un \emph{fork} del repositorio
\url{https://github.com/andres-karoll/personapp-hexa-spring-boot},
trabajar sobre la copia del grupo y conservar un historial Git. No deben
presentar como entrega un vínculo al repositorio original del profesor.

\textbf{Resultado (HU-001):} el repositorio del equipo es
\url{https://github.com/SarAlbN1/AS-PersonAPP_Laboratorio2}, fork de la plantilla.
Los cambios se integran a \texttt{main} mediante una rama por historia y un PR
revisado por otro integrante (PR \#30, 2026-11-08).

\subsubsection*{Actividad 2. Preparar el ambiente}
La ejecución inicial se realiza respetando las versiones de la plantilla:
Java 11, Spring Boot 2.7.11, las dependencias Maven y los motores de base
de datos. Se debe habilitar Lombok en el IDE si es necesario y comprobar
la configuración de los dos adaptadores de entrada. Como mínimo,
registrar las salidas de \texttt{java -version}, \texttt{mvn -version} y
la versión de cada motor de base de datos.

\textbf{Resultado (HU-003):} Temurin 11.0.32.1, Maven 3.9.9, MariaDB 11.4.13 y
MongoDB 6.0.28 en contenedores Docker. Se eligió MongoDB 6.0 por ser la versión
más alta que admite el controlador Java 4.6 incluido en Spring Boot 2.7.11
(PR \#33, 2026-11-08).

\subsubsection*{Actividad 3. Ejecutar los scripts DDL/DML}
La carpeta \texttt{scripts} del proyecto contiene los archivos de MariaDB
y MongoDB. La ejecución debe crear estructuras y cargar registros de prueba.
Se detecta un punto de revisión en los archivos base: la configuración
Spring señala \texttt{persona\_db} para MongoDB, mientras que el
\texttt{persona\_dml\_mongo.js} selecciona \texttt{prueba\_db}.
Antes de probar, el grupo debe verificar y unificar la base de datos
utilizada. Las contraseñas de ejemplo son propias de un ambiente local de
laboratorio y no deben asumirse como configuración segura de producción.

\textbf{Resultado (HU-004):} los cuatro scripts se ejecutaron en contenedores
limpios. El DML de MongoDB se corrigió para cargar en \texttt{persona\_db}, la base
que usan REST y CLI. Ambos motores contienen las mismas cinco personas
(PR \#34, 2026-11-09).

\subsubsection*{Actividad 4. Compilar y ejecutar la versión base}
Desde el directorio raíz Maven se puede intentar la compilación con:
\begin{verbatim}
mvn clean install
\end{verbatim}
El proyecto contiene dos clases de arranque independientes:
\texttt{PersonAppRestApi} y \texttt{PersonAppCli}. Se deben ejecutar
ambos clientes, verificar el endpoint inicial de personas y comprobar
el acceso a Swagger y al menú CLI.

\textbf{Resultado (HU-005):} la compilación termina con \texttt{BUILD SUCCESS}.
Los JAR generados no son ejecutables, porque \texttt{spring-boot-maven-plugin}
no declara \texttt{repackage}, y \texttt{spring-boot:run} falla dentro de cada
módulo por la propiedad \texttt{\$\{revision\}}. Por eso ambas aplicaciones se
arrancaron con el classpath que calcula el reactor Maven, sin cambiar el código.
Los resultados están en la tabla de pruebas ejecutadas (PR \#35, 2026-11-09).
Los JAR ejecutables y la resolución de \texttt{\$\{revision\}} se corrigen al
actualizar los POM en la fase 2.

\subsubsection*{Actividad 5. Identificar y corregir el Mapper}
El profesor advierte expresamente que hay un error en un Mapper de la
plantilla. Su resolución debe documentarse como una actividad de
\emph{depuración}: reproducir el error, registrar su mensaje, localizar
la conversión defectuosa, explicar la causa, aplicar la corrección y
volver a ejecutar la prueba.

\textbf{Resultado (HU-006):} el \texttt{POST /api/v1/persona} respondía 500 en
ambos motores con \texttt{NullPointerException: gender is marked non-null but is
null} en \texttt{validateGenero} de los mappers de salida. La causa era
\texttt{PersonaMapperRest.fromAdapterToDomain}, que estaba sin implementar y
devolvía un \texttt{Person} vacío; los mappers de MariaDB y MongoDB solo eran los
primeros en recibirlo. La corrección se limitó a ese método: convierte
\texttt{dni} y \texttt{age} a enteros, copia los nombres y traduce el sexo
(\texttt{M}/\texttt{MALE} y \texttt{F}/\texttt{FEMALE}). Seis pruebas de regresión
cubren la transformación, y después el alta guarda en los dos motores (PR \#38,
2026-11-11).

\subsection{Fase 2. Actualizar las tecnologías}
Criterio de salida: todos los módulos compilan con las versiones finales, REST y
CLI arrancan, ambos motores responden y las pruebas de regresión de los mappers
siguen pasando.

El profesor aclaró que actualizar una plataforma no es solo cambiar el número de
versión en el código: hay que resolver las dependencias de una versión a otra y
los cambios de las API que usa el proyecto. HU-007 registra las versiones
iniciales, adopta la LTS más alta compatible para Java y para cada tecnología que
publique versiones LTS, y elige para las demás la versión estable más reciente con
soporte vigente. Compara el árbol de dependencias antes y después, revisa los POM,
plugins, anotaciones, API y configuración, y añade la ejecución
\texttt{repackage} en REST y CLI y \texttt{flatten-maven-plugin} en el POM padre.

\textbf{Pendiente del equipo (HU-007):} versiones iniciales y finales con su
fuente, cambios de dependencias y de API, errores encontrados y su solución, y
resultados de compilación, arranque y pruebas.

\subsection{Fase 3. Contenedorizar la aplicación}
Criterio de salida: desde un clon limpio se construyen las imágenes, se levantan
REST, CLI, MariaDB y MongoDB, y la aplicación actualizada funciona con ambos
motores.

HU-034 versiona los Dockerfile de las entradas REST y CLI y un Docker Compose con
ambos motores. Externaliza la configuración y las credenciales, define volúmenes
persistentes y documenta la carga de los scripts DDL/DML. Los contenedores
empaquetan las entradas y las bases existentes; ningún módulo Maven se convierte
en un servicio aparte.

\textbf{Pendiente del equipo (HU-034):} imágenes y Compose versionados,
comandos de construcción y arranque, persistencia tras recrear los contenedores
y procedimiento de parada y limpieza.

\subsection{Fase 4. Completar las funcionalidades}
Criterio de salida: Persona, Profesión, Teléfono y Estudios ofrecen las seis
operaciones acordadas por el equipo (consultar por identificador, listar todos,
contar, crear, editar y eliminar), el listado parcial y la consulta de una
persona con sus teléfonos, en REST y CLI, con persistencia en ambos motores.

Según el profesor, la funcionalidad de la plantilla se concentra en Persona.
El equipo completará primero Persona de punta a punta: conteo y listado parcial
en los puertos y adaptadores de salida, y rutas REST y opciones CLI para las seis
operaciones. Después seguirán Profesión, Teléfono y Estudios. Para estos recursos
ya existen modelos, entidades, documentos y mappers. Faltan los puertos, casos de
uso, repositorios, adaptadores de salida y entradas REST y CLI. El listado parcial
usa páginas con orden estable por clave. El conteo devuelve el total del motor
elegido, y cero si la colección está vacía.

\textbf{Pendiente del equipo:} clases creadas o modificadas,
justificación de las decisiones, endpoints finales, menús y pruebas.

\subsection{Verificación y entrega}
Criterio de salida: hay evidencias de las seis operaciones en ambos clientes y
motores sobre el despliegue en contenedores, Swagger 3 documenta la API final, el
README es reproducible, el informe está completo y existe la URL del tag.

\subsubsection*{Pruebas funcionales y de integración}
Se comprobarán las seis operaciones de cada recurso desde REST y CLI con
MariaDB y MongoDB, con las versiones finales y en los contenedores. Las evidencias deben incluir listados con orden estable,
conteo antes y después de un alta o baja y aislamiento entre motores. Por
ejemplo, con tres registros y tamaño de página dos, la primera página debe
devolver dos registros y la segunda uno, sin duplicados; el conteo debe pasar de
tres a dos tras una eliminación, sin alterar el otro motor.

Además de las operaciones exitosas, se documentarán entradas inválidas,
identificadores inexistentes, relaciones entre entidades, y el cambio
\textbf{en caliente} entre bases durante la misma ejecución. En REST
se registrarán método, ruta, cuerpo de solicitud, respuesta y código HTTP;
en CLI, opción seleccionada y salida observada. Deben conservarse
capturas útiles y consistentes con la versión etiquetada.

\subsubsection*{README, publicación y etiqueta final}
La entrega técnica deberá incluir un README que explique requisitos,
configuración del ambiente, preparación de MariaDB/MongoDB, ejecución de
scripts, compilación Maven, arranque de REST y CLI, uso de Swagger y
comandos de prueba. También debe contener los scripts DDL/DML y el
código fuente, y publicarse con un tag Git que identifique el estado
final entregado.

\begin{table}[H]
\centering\small
\begin{tabularx}{\textwidth}{@{}p{3.8cm}X@{}}
\toprule
\textbf{Entregable}&\textbf{Estado al 2026-11-11}\\\midrule
README reproducible&\texttt{README.md} declara el stack y los pasos de preparación, scripts y arranque. Falta comprobarlo desde un clon limpio.\\
Scripts DDL/DML&\path{scripts/} contiene \path{persona_ddl_maria.sql}, \path{persona_dml_maria.sql}, \path{persona_ddl_mongo.js} y \path{persona_dml_mongo.js}. HU-004 unificó MongoDB en \texttt{persona\_db}.\\
Código fuente&Siete módulos Maven. El alta de Persona funciona desde HU-006 (PR \#38). Persona sigue parcialmente expuesta y los otros tres recursos no tienen entradas.\\
URL del tag&No existe tag en el repositorio. \textbf{Pendiente.}\\
\bottomrule
\end{tabularx}
\caption{Cobertura de los entregables técnicos.}
\end{table}

\textbf{Pendiente del equipo:} URL definitiva del repositorio,
URL del tag, identificación de commit, resultados de pruebas y
confirmación de que el README reproduce el despliegue.

\section{Conclusiones y lecciones aprendidas}

\subsection{Conclusiones preliminares de diseño}
La organización por puertos, casos de uso y adaptadores proporciona una
base para separar las responsabilidades del negocio de REST, CLI y de
los mecanismos de persistencia. En comparación con la estructura usada
para el Laboratorio 1, este diseño permite estudiar explícitamente la
intercambiabilidad de adaptadores y el papel de las dependencias
entre módulos.

La plantilla permite identificar dos entradas y dos tecnologías de
persistencia, así como contratos de aplicación que deben completarse y
exponerse. La validación de la actualización tecnológica, del despliegue en
contenedores y de las seis operaciones queda sujeta a la implementación y a las
pruebas reales del equipo.

\subsection{Lecciones aprendidas por documentar}
\textbf{Pendiente del equipo:} describir dificultades reales en la
compilación y ejecución inicial, el diagnóstico del Mapper, la
preparación de los motores, la actualización de versiones y sus dependencias,
la contenedorización, la implementación de los dominios
faltantes y la selección dinámica de MariaDB/MongoDB. Esta subsección
se completará al terminar las actividades para no atribuir al grupo
experiencias que todavía no se han documentado.

\clearpage
\begin{thebibliography}{9}
\bibitem{plantilla} Sánchez Martín, A. A. (s. f.). \textit{PersonAPP: plantilla de arquitectura limpia con Spring Boot}.
Repositorio de GitHub. \url{https://github.com/andres-karoll/personapp-hexa-spring-boot}.
\bibitem{cockburn} Cockburn, A. (2005). \textit{Hexagonal Architecture}.
\url{https://alistair.cockburn.us/hexagonal-architecture/}.
\bibitem{martin} Martin, R. C. (2017). \textit{Clean Architecture: A Craftsman's Guide to Software Structure and Design}. Prentice Hall.
\bibitem{maven} Apache Software Foundation. (s. f.). \textit{Maven Guides}.
\url{https://maven.apache.org/guides/}.
\bibitem{spring} Spring. (s. f.). \textit{Spring Boot reference documentation}.
\url{https://docs.spring.io/spring-boot/}.
\bibitem{fowler} Fowler, M. (2011). \textit{CQRS}.
\url{https://martinfowler.com/bliki/CQRS.html}.
\bibitem{mariadb} MariaDB Foundation. (s. f.). \textit{MariaDB documentation}.
\url{https://mariadb.com/docs/}.
\bibitem{mongodb} MongoDB. (s. f.). \textit{MongoDB manual}.
\url{https://www.mongodb.com/docs/manual/}.
\bibitem{openapi} OpenAPI Initiative. (s. f.). \textit{OpenAPI Specification}.
\url{https://spec.openapis.org/oas/latest.html}.
\end{thebibliography}
\end{document}
